\documentclass[11pt]{article}
\usepackage[margin=0.75in]{geometry}
\usepackage{microtype}
\usepackage{parskip}
\usepackage{helvet}
\usepackage{array}
\usepackage{xcolor}
\usepackage{booktabs}
\usepackage{multicol}
\renewcommand{\familydefault}{\sfdefault}

\definecolor{gwblue}{RGB}{0,59,92}
\definecolor{gwbuff}{RGB}{168,153,104}

% No page numbers
\pagestyle{empty}

\begin{document}

%% ---- Header bar ----
\noindent{\large\bfseries\color{gwblue}
GWSPH Research Showcase \hfill Biostatistics \& Bioinformatics}

\vspace{0.4em}
\noindent\rule{\linewidth}{1.5pt}
\vspace{0.2em}

%% ---- Title ----
\noindent{\large\bfseries
Computational Pipeline Choice Drives Shotgun-Amplicon\\
Discordance in Microbial Community Profiling}

\vspace{0.8em}

%% ---- Two-column body: abstract left, sidebar right ----
\noindent
\begin{minipage}[t]{0.67\linewidth}
\raggedright

\noindent\textbf{Background \& Significance:} Reported discordance between
shotgun and amplicon microbiome profiling is often attributed to biological
or technical wet-lab differences, but the contribution of computational
pipelines remains underexplored. Pipeline choice is a potentially modifiable
source of bias with direct implications for reproducibility in microbiome
research.

\medskip

\noindent\textbf{Purpose:} We ask whether routing shotgun data through the
amplicon computational pipeline recovers concordant community profiles, and
whether Fisher-Rao distance provides more reliable comparisons than
Bray-Curtis dissimilarity across pipeline streams.

\medskip

\noindent\textbf{Methods:} Using data from a published vermicomposting
experiment across 12 microcosms in a $2{\times}2$ factorial design
(earthworm presence $\times$ sampling time), we compared three analytical
streams: native amplicon 16S sequencing, contig-based 16S extraction from
shotgun assemblies, and shotgun reads processed through the amplicon DADA2
pipeline. Community dissimilarities were computed using Fisher-Rao distance
and Bray-Curtis dissimilarity, followed by PERMANOVA\@.

\medskip

\noindent\textbf{Findings \& Conclusions:} Contig-based profiles showed
strong concordance with native amplicon results across all taxonomic levels,
while read-based profiles diverged substantially. Fisher-Rao distance
produced more consistent pairwise estimates with lower variance, yielding
more reliable separation among pipeline streams than Bray-Curtis.
These results suggest that much reported shotgun-amplicon discordance
reflects dry-lab pipeline differences rather than biological signal.

\medskip

\noindent\textbf{Implications:} Future work will extend this framework to
fungal analysis and validation against reference mock communities to assess
implications for clinical and public health microbiome applications.

\end{minipage}
\hfill
\begin{minipage}[t]{0.28\linewidth}
\raggedright
\setlength{\parskip}{0.6em}

\noindent\textcolor{gwblue}{\rule{\linewidth}{1pt}}

\noindent{\small\bfseries\color{gwblue} Primary Presenter}

\noindent{\small Clark Gaylord}

\noindent\textcolor{gwblue}{\rule{\linewidth}{0.5pt}}

\noindent{\small\bfseries\color{gwblue} Co-Presenter(s)}

\noindent{\small N/A}

\noindent\textcolor{gwblue}{\rule{\linewidth}{0.5pt}}

\noindent{\small\bfseries\color{gwblue} Status}

\noindent{\small Doctoral}

\noindent\textcolor{gwblue}{\rule{\linewidth}{0.5pt}}

\noindent{\small\bfseries\color{gwblue} Authors}

\noindent{\small Clark Gaylord\\
Marcos P\'{e}rez-Losada\\
Keith A.\ Crandall}

\noindent\textcolor{gwblue}{\rule{\linewidth}{0.5pt}}

\noindent{\small\bfseries\color{gwblue} Research Mentor /\\Department Chair}

\noindent{\small Keith A.\ Crandall}

\noindent\textcolor{gwblue}{\rule{\linewidth}{0.5pt}}

\end{minipage}

\end{document}
